\section{Evaluation Scope and Data Splits}
\label{sec:protocol}

Section~\ref{sec:method} defines the calibrated detector and embeds the no-leakage data roles used by the method. This section specifies how those roles are instantiated for the main evidence chain and how supplementary datasets are interpreted.

\subsection{Main Evidence Chain}

The main experimental evidence is based on CICIDS2017 under the method-level data separation defined in Section~\ref{sec:method}. CICIDS2017 is used as the primary dataset because it provides temporally organized flow records, multiple attack types, and a common benchmark setting for intrusion-detection research. Main claims about calibrated low-noise detection, threshold behavior, TAD score ablation, attack-type breakdown, and auditable alarm evidence are drawn from this CICIDS2017 setting.

\subsection{CICIDS2017 Split}
\label{sec:cicids_split}

For the CICIDS2017 main evaluation, we use a day-separated split rather than randomly mixing windows across the full dataset. Traffic from earlier weekdays is used on the training and calibration side, while Friday traffic is reserved for final testing. This preserves temporal separation between detector construction and final evaluation and reduces the risk that adjacent or near-duplicate traffic windows leak across train and test partitions.

Although the training-side raw files contain both benign and attack records, only benign windows are used for scaler fitting, model training, TAD reference construction, and independent threshold calibration. Attack windows from the training-side files are excluded from these steps and are not used to tune the detector. Test labels are used only after threshold fixation, for reporting detection metrics and attack-type breakdowns.

The exact CICIDS2017 file list used for this split is provided in the artifact documentation to support reproducibility. In the main paper, we report the protocol-level settings used throughout the main evidence chain: window size 128, stride 16, anomaly-ratio threshold 0.15, and independent benign calibration ratio 0.2.

\subsection{Supplementary Dataset Scope}

SWaT and UNSW-NB15 are included only as supplementary operating-boundary checks. SWaT provides an industrial-control validation context, while UNSW-NB15 is used as a distribution-shift stress test. These datasets are not used as co-equal primary benchmarks and are not presented as comprehensive cross-dataset generalization evidence. In particular, low UNSW-NB15 recall is treated as threshold-transfer failure under distribution shift, not as a success case.
